\documentclass[10pt,a4paper]{amsart}
\usepackage[margin=19mm]{geometry}
\usepackage{amsmath,amssymb,mathtools}
\usepackage[scr=rsfs,cal=euler]{mathalfa}
\usepackage{xfrac}
\usepackage[unicode,hidelinks,bookmarks=false]{hyperref}
\newcommand{\ie}{i.e.\ }
\newcommand{\cf}{cf.\ }
\newcommand{\resp}{respectively\ }
\newcommand{\Subsection}{\subsection}
\providecommand{\nobreakdash}{\nobreak\hbox{-}}
\setlength{\emergencystretch}{2em}
\multlinegap0pt
\usepackage{amssymb}
\usepackage{xcolor}
\newtheorem{lemma}{Lemma}
\newcommand\fdem{$\Box$}
\title{On the integro-differential equation arising in the ruin problem for annuity payment models}
\author{Platon Promyslov}
\date{Source: arXiv:2601.01447v1 (CC BY 4.0)}
\begin{document}
\maketitle

\noindent\emph{Source and licence.} On the integro-differential equation arising in the ruin problem for annuity payment models. Version arXiv:2601.01447v1 (CC BY 4.0), distributed by arXiv under the Creative Commons Attribution 4.0 International licence, http://creativecommons.org/licenses/by/4.0/, as recorded in the arXiv licence metadata for this exact version and shown on the abstract page https://arxiv.org/abs/2601.01447v1. Full licence notice: \texttt{licenses/arxiv-2601.01447-CC-BY-4.0.txt}.\par\smallskip

\noindent\textit{Editorial scope.} Counted unit: the article's lemma lemma (source0.tex lines 338--340). The article's complete original proof of it is delivered with it (source0.tex lines 341--360, one \texttt{proof} environment of 20 lines). No mathematics is written, repaired or re-derived: the statement and the proof are the article's own lines, in the article's own notation, and no retained line is substituted or re-flowed.  No further source-local context is needed: every cross-reference of every retained line resolves inside the two delivered spans.  Nothing was omitted from any retained span, so the package carries no omission record.  No retained line contains a citation, so the wrapper carries no bibliography and no bibliography entry.  No retained line contains a figure environment, an image-inclusion command or drawing code, so no image is delivered and the package declares no asset dependency.  Editorial formatting only, in addition: an AMS \texttt{amsart} preamble that restates the article's own declarations and macros used by the retained lines, namely the environments \texttt{lemma} on separate counters, together with the standard packages those lines need.  The retained environments print in this document as Lemma 1 in order of appearance, whereas the article prints the same items with the numbering of its own full document.  The complete native source of the article as distributed in the arXiv e-print for this exact version is bundled unchanged as \texttt{sources/192080/source0.tex}.
\begin{lemma}
$\Phi(0+)=\Phi(0):=0$.
\end{lemma}

\begin{proof} 
Consider the stochastic exponential $Z=(Z_t)_{t\ge 0}$ given by the linear SDE 
$$
dZ_t=Z_t \left( ((a-r)\kappa + r) \,d t + \sigma \,dW_t \right), \qquad Z_0=1, 
$$
and define the process $\hat X^u$ 
$$
\hat X^u_t=Z_t\Big( u-c\int_0^t Z_s^{-1}ds\Big).
$$

Let $\hat \tau^ u:=\inf\{t\ge 0\colon \hat X^u=0\}$ and let $T_1$ be the instant of first jump of the Poisson process $N$. 
According to the stochastic Cauchy formula giving the solution of non-homogeneous linear SDE 
$X^u=\hat X^u$ on $[0,T_1)$ and $X^u_{T_1}=\hat X^u_{T_1}+\Delta X^u_{T_1}\ge \hat X^u_{T_1}$. Thus, $\hat \tau^ u\wedge T_1= \tau^ u\wedge T_1$. It follows that 
$$
0\le \hat X^u_{\tau^u\wedge T_1}=Z_{\tau^u\wedge T_1}\Big( u-c\int_0^{\tau^u\wedge T_1} Z_s^{-1}ds\Big)
\le \xi^* u-c\frac {\xi_*}{\xi^*}(\tau^u\wedge T_1),
$$
where $\xi_*:=\inf_{s\le T_1}Z_s$, $\xi^*:=\sup_{s\le T_1}Z_s$. Therefore, $\tau^ u\to 0$ as $u\to 0$ and the result follows.
\fdem
\end{proof} 

\end{document}
